%-------------------------------------------------------------------------------
% CV per il bando "Tutorato alla pari DISI 2026/27" (Allegato 8), redatto secondo
% lo schema "CV_Modello_Europeo" allegato al bando.
%
% I dati personali richiesti dallo schema non sono nel repository pubblico:
%   - telefono:  cv/private-contact.tex          (\PrivatePhone)
%   - indirizzo, cittadinanza, data di nascita, sesso:
%                cv/europass/private/dati-personali.tex
% Entrambi i file sono ignorati da git. Compilare in locale con:
%   cd cv/europass && latexmk -pdf CV-Tutorato-DISI-2026.tex
%-------------------------------------------------------------------------------
\documentclass[10pt,a4paper]{article}

\usepackage[utf8]{inputenc}
\usepackage[T1]{fontenc}
\usepackage[italian]{babel}
\usepackage{helvet}
\renewcommand{\familydefault}{\sfdefault}
\usepackage[margin=1.6cm, top=1.5cm, bottom=1.6cm]{geometry}
\usepackage{array}
\usepackage{xcolor}
\usepackage{enumitem}
\usepackage[hidelinks]{hyperref}
\urlstyle{same}
\pagestyle{plain}
\setlength{\parindent}{0pt}
\setlist[itemize]{nosep, leftmargin=1em, label=--}

\definecolor{ecvblue}{RGB}{14, 65, 148}
\definecolor{ecvgray}{RGB}{90, 90, 90}

% Private data (see header)
\IfFileExists{../private-contact.tex}{\input{../private-contact.tex}}{\newcommand{\PrivatePhone}{\textcolor{red}{DA COMPILARE}}}
\IfFileExists{private/dati-personali.tex}{\input{private/dati-personali.tex}}{%
  \newcommand{\DatiIndirizzo}{\textcolor{red}{DA COMPILARE}}%
  \newcommand{\DatiCittadinanza}{\textcolor{red}{DA COMPILARE}}%
  \newcommand{\DatiNascita}{\textcolor{red}{DA COMPILARE}}%
  \newcommand{\DatiSesso}{\textcolor{red}{DA COMPILARE}}}

% Two-column Europass layout: label on the left, content on the right
\newcolumntype{L}{>{\raggedleft\arraybackslash}p{4.6cm}}
\newcolumntype{R}{>{\raggedright\arraybackslash}p{12.3cm}}
\newcommand{\ecvrow}[2]{%
  \noindent\begin{tabular}{@{}L!{\color{ecvblue}\vrule width 0.6pt}R@{}}\leavevmode{\small\color{ecvgray}#1} & #2\end{tabular}\par\vspace{2pt}}
\newcommand{\ecvsec}[1]{%
  \vspace{8pt}\noindent\begin{tabular}{@{}L!{\color{ecvblue}\vrule width 0.6pt}R@{}}%
  \leavevmode{\bfseries\color{ecvblue}\MakeUppercase{#1}} & \end{tabular}\par\vspace{2pt}}
\newcommand{\ecvgap}{\vspace{6pt}}
\newcommand{\lev}[1]{\textbf{#1} \small Utente avanzato}

\begin{document}

\ecvrow{}{{\Large\bfseries\color{ecvblue} Curriculum Vitae Europass}}
\vspace{6pt}

%-------------------------------------------------------------------------------
\ecvsec{Informazioni personali}
\ecvrow{Nome(i) / Cognome(i)}{\textbf{\large Mattia Giovannini}}
\ecvrow{Indirizzo(i)}{\DatiIndirizzo}
\ecvrow{Telefono(i)}{Cellulare: \PrivatePhone}
\ecvrow{E-mail}{\href{mailto:mattia.giovannini6@studio.unibo.it}{mattia.giovannini6@studio.unibo.it}}
\ecvrow{Cittadinanza}{\DatiCittadinanza}
\ecvrow{Data di nascita}{\DatiNascita}
\ecvrow{Sesso}{\DatiSesso}

%-------------------------------------------------------------------------------
\ecvsec{Esperienza di tirocinio o professionale}

\ecvrow{Date}{2021 -- oggi}
\ecvrow{Lavoro o posizione ricoperti}{\textbf{Tutor privato (ripetizioni) di Matematica, Fisica e Informatica}}
\ecvrow{Principali attività e responsabilità}{%
  \begin{itemize}
    \item Lezioni individuali per studenti delle scuole superiori di ogni indirizzo, in presenza a Bologna e online
    \item Recupero di debiti formativi, preparazione di verifiche ed esami e supporto continuativo durante l'anno scolastico
    \item Percorsi personalizzati a partire dalle lacune del singolo studente, con l'obiettivo di renderlo autonomo nello studio
  \end{itemize}}
\ecvrow{Nome e indirizzo del datore di lavoro}{Attività in proprio, Bologna e online}
\ecvrow{Tipo di attività o settore}{Istruzione -- supporto allo studio}
\ecvgap

\ecvrow{Date}{2016 -- 2020}
\ecvrow{Lavoro o posizione ricoperti}{\textbf{Educatore volontario presso centro estivo}}
\ecvrow{Principali attività e responsabilità}{Animazione e gestione di gruppi di bambini e ragazzi in attività educative, lavoro di squadra con gli altri educatori}
\ecvrow{Nome e indirizzo del datore di lavoro}{Estate Ragazzi, Bologna}
\ecvrow{Tipo di attività o settore}{Attività educative estive (volontariato)}
\ecvgap

\ecvrow{Date}{2020 -- 2021}
\ecvrow{Lavoro o posizione ricoperti}{\textbf{Webmaster ed editor studentesco}}
\ecvrow{Principali attività e responsabilità}{Progettazione e gestione delle piattaforme studentesche dell'istituto (\textit{Righi Network} e \textit{Il Sottobanco}) e della presenza online del giornale scolastico}
\ecvrow{Nome e indirizzo del datore di lavoro}{Liceo Scientifico Augusto Righi, Bologna}
\ecvrow{Tipo di attività o settore}{Scuola -- attività studentesca (volontariato)}

%-------------------------------------------------------------------------------
\ecvsec{Istruzione e formazione}

\ecvrow{Date}{2026 -- 2028 (in corso)}
\ecvrow{Titolo della qualifica rilasciata}{\textbf{Laurea Magistrale in Ingegneria Informatica} (doppio titolo con il Master in Ubiquitous Systems and Networks)}
\ecvrow{Principali tematiche/competenze professionali acquisite}{Ingegneria informatica, sistemi distribuiti e ubiqui, reti, piattaforme enterprise e DevOps, ottimizzazione}
\ecvrow{Nome e tipo d'organizzazione erogatrice dell'istruzione e formazione}{Università di Bologna (università pubblica) \& Université Côte d'Azur -- Polytech Nice Sophia, Nizza (Francia)}
\ecvrow{Livello nella classificazione nazionale o internazionale}{Laurea Magistrale -- EQF 7}
\ecvgap

\ecvrow{Date}{2021 -- luglio 2026}
\ecvrow{Titolo della qualifica rilasciata}{\textbf{Laurea in Ingegneria Informatica} -- voto finale 94/110}
\ecvrow{Principali tematiche/competenze professionali acquisite}{%
  Analisi Matematica, Algebra Lineare e Geometria, Probabilità e Statistica, Fondamenti di Informatica, Sistemi Operativi, Reti di Calcolatori, Sistemi Informativi, Ingegneria del Software, Tecnologie Web, Architettura degli Elaboratori.\newline
  Tesi: \textit{Rilevamento del Disturbo dello Spettro Autistico mediante Machine Learning applicato a Dati di Eye Tracking}}
\ecvrow{Nome e tipo d'organizzazione erogatrice dell'istruzione e formazione}{Università di Bologna (università pubblica)}
\ecvrow{Livello nella classificazione nazionale o internazionale}{Laurea -- EQF 6}
\ecvgap

\ecvrow{Date}{2025}
\ecvrow{Titolo della qualifica rilasciata}{\textbf{Mobilità Erasmus+ per studio}}
\ecvrow{Principali tematiche/competenze professionali acquisite}{Progetto di machine learning per l'analisi del disturbo dello spettro autistico su dati di eye-tracking; studio e lavoro in un contesto accademico internazionale e francofono}
\ecvrow{Nome e tipo d'organizzazione erogatrice dell'istruzione e formazione}{Université de Tours, Tours (Francia) -- università pubblica}
\ecvrow{Livello nella classificazione nazionale o internazionale}{Periodo di studio all'estero durante la Laurea}
\ecvgap

\ecvrow{Date}{2016 -- 2021}
\ecvrow{Titolo della qualifica rilasciata}{\textbf{Diploma di Liceo Scientifico}}
\ecvrow{Nome e tipo d'organizzazione erogatrice dell'istruzione e formazione}{Liceo Scientifico Augusto Righi, Bologna (scuola secondaria di secondo grado)}
\ecvrow{Livello nella classificazione nazionale o internazionale}{Diploma di istruzione secondaria superiore -- EQF 4}

%-------------------------------------------------------------------------------
\ecvsec{Capacità e competenze personali}

\ecvrow{Madrelingua(e)}{\textbf{Italiano}}
\ecvgap
\ecvrow{Altra(e) lingua(e)}{}
\ecvrow{Autovalutazione}{%
  {\footnotesize
  \begin{tabular}{@{}l|cc|cc|c@{}}
    & \multicolumn{2}{c|}{\textbf{Comprensione}} & \multicolumn{2}{c|}{\textbf{Parlato}} & \textbf{Scritto} \\
    \textit{Livello europeo (*)} & Ascolto & Lettura & Interazione orale & Produzione orale & \\
    \hline
    \textbf{Inglese} & C1 & C1 & C1 & C1 & C1 \\
    \textbf{Francese} & C1 & C1 & C1 & C1 & C1 \\
  \end{tabular}}\par\vspace{3pt}
  {\footnotesize Inglese: CLA C1, Università di Bologna; Duolingo English Test 130/160 (C1); Cambridge First Certificate B2 (178/190).\newline
  Francese: CLA C1, Università di Bologna; corsi B2/C1 presso l'Université de Tours.\newline
  (*) Quadro Comune Europeo di Riferimento per le Lingue -- C1: utente avanzato.}}
\ecvgap

\ecvrow{Capacità e competenze sociali}{%
  \begin{itemize}
    \item Insegnamento individuale a studenti di età, indirizzi e livelli diversi, acquisito in cinque anni di ripetizioni: capacità di ascolto, di spiegare in modo chiaro e di adattare il metodo allo studente
    \item Gestione di gruppi di bambini e ragazzi come educatore volontario a Estate Ragazzi (2016--2020)
    \item Vita e studio in un ambiente multiculturale durante il semestre Erasmus+ a Tours
    \item Lavoro di squadra in progetti universitari multidisciplinari (BOOM GenAI Challenge 2026, Lean Startup Lab)
  \end{itemize}}
\ecvrow{Capacità e competenze organizzative}{%
  \begin{itemize}
    \item Organizzazione autonoma dell'attività di ripetizioni: pianificazione delle lezioni e dei percorsi di recupero, gestione del calendario con più studenti
    \item Coordinamento delle piattaforme studentesche e della redazione online del giornale scolastico al Liceo Righi (2020--2021)
    \item Gestione di progetti di gruppo universitari con scadenze e ruoli condivisi
  \end{itemize}}
\ecvrow{Capacità e competenze informatiche}{%
  \begin{itemize}
    \item Linguaggi: C, Java, Python, JavaScript, HTML/CSS, SQL, Bash
    \item Machine learning e dati: scikit-learn, NumPy, pandas, Matplotlib, Jupyter Notebook
    \item Sistemi e reti: Linux, TCP/IP, Docker, Podman, Kubernetes, Git, \LaTeX, MATLAB
  \end{itemize}
  Acquisite durante la Laurea e la Laurea Magistrale in Ingegneria Informatica e in progetti personali (homelab self-hosted).}

%-------------------------------------------------------------------------------
\ecvsec{Ulteriori informazioni}
\ecvrow{Progetti}{%
  \begin{itemize}
    \item \textbf{Scheduler Kubernetes resource-aware} (2026, in corso) -- project work, Enterprise Platforms and DevOps M
    \item \textbf{Benchmark di ripristino da guasti di link SDN} (2026, in corso) -- Discrete and Network Optimization Problems M
    \item \textbf{AgenticOS}, audit multi-agente con IA (2026) -- Lean Startup Lab, Technology Entrepreneurship
    \item \textbf{BOOM GEN AI Challenge 2026} -- assistente GenAI per l'analisi di bandi di gara, con VEM Sistemi
  \end{itemize}}
\ecvrow{Portfolio}{\href{https://mattia-giovannini.github.io}{mattia-giovannini.github.io} \quad \href{https://github.com/frigori}{github.com/frigori}}

%-------------------------------------------------------------------------------
\ecvsec{Allegati}
\ecvrow{}{Copia fronte/retro di un documento di identità valido (allegata alla domanda)}

%-------------------------------------------------------------------------------
\ecvsec{Firma}
\ecvrow{}{Autorizzo il trattamento dei miei dati personali ai sensi del Decreto Legislativo 30 giugno 2003, n. 196 ``Codice in materia di protezione dei dati personali''.\par\vspace{14pt}
  Bologna, \today \hfill Mattia Giovannini\hspace*{1.5cm}}

\end{document}
